\documentclass[12pt,a4paper]{article}
\usepackage[utf8]{inputenc}
\usepackage[T1]{fontenc}
\usepackage{geometry}
\usepackage{booktabs}
\usepackage{amsmath}
\usepackage{hyperref}
\usepackage{setspace}
\usepackage{titlesec}
\usepackage{enumitem}
\usepackage{array}
\usepackage{xcolor}

\geometry{margin=2.5cm}
\onehalfspacing

\hypersetup{
    colorlinks=true,
    linkcolor=blue,
    urlcolor=blue,
    citecolor=blue
}

\title{Statistical Analysis of the Phaistos Disc:\\
A Computational Methodology for Phonetic Key Evaluation}

\author{Manolis Chavadakis\\
\small Independent Researcher\\
\small \href{mailto:aikopappas@gmail.com}{aikopappas@gmail.com}}

\date{June 2026 \quad Version 3.3}

\begin{document}

\maketitle

\begin{abstract}
The Phaistos Disc ($\sim$1700 BCE) remains one of archaeology's most debated undeciphered objects. We present a blind computational framework for evaluating competing phonetic key hypotheses, applying Bonferroni-corrected Monte Carlo simulation across 10 candidate keys scored against three reference corpora: Luwian Hieroglyphic vocabulary, Linear A frequency tables, and the AED-TEI Egyptian corpus (675,773 tokens from 13,950 texts). Three key-independent findings are established irrespective of any phonetic assumption: (1)~the [Sign\#36$\to$Sign\#11] sequential bigram shows $Z=10$ excess adjacency (obs/exp$=7.69\times$, $p\approx 0$); (2)~corpus-domain control confirms ritual text classification (theological $Z=27.16$ vs.\ administrative $Z=-0.40$); (3)~Sign~\#45 (solar rosette) appears exclusively at spiral centers A31 and B30. The Luwian Hieroglyphic key (\textsc{G\_Luwian}) achieves the highest Bonferroni-significant score ($S=523$, $p<0.0001$), yielding a solar-water cosmological reading structurally parallel to the Egyptian \textit{Amduat}. A negative control test on a synthetic disc ($Z=1.99$, not significant) establishes that token-level scores are $\sim$94\% frequency-driven. Accordingly, only the three key-independent pillars are presented as primary claims. All code and data are released open-source.

\medskip
\noindent\textbf{Keywords:} Phaistos Disc, undeciphered scripts, computational linguistics, Luwian hieroglyphics, Monte Carlo simulation, Bonferroni correction, Bronze Age Aegean, ritual text analysis
\end{abstract}

\tableofcontents
\newpage

%------------------------------------------------------------
\section{Introduction}
%------------------------------------------------------------

The Phaistos Disc, discovered in 1908 at the Minoan palace of Phaistos (Crete) and dated to approximately 1700~BCE, bears 241 impressed signs from a repertoire of 45 distinct symbols arranged in a double-sided spiral across 61 word-groups. It remains unique: no second exemplar exists, and its script, language, and reading direction have not been established to scholarly consensus.

Previous decipherment attempts number in the hundreds and span proposed languages from Minoan to Phoenician, Greek, Anatolian, and Semitic. Nearly all share a methodological weakness: the proposed phonetic key is constructed to produce semantically plausible readings, creating unfalsifiable circularity.

The present study does not propose a decipherment. It proposes a \textbf{statistical methodology} for ranking competing phonetic key hypotheses against objective reference corpora, with explicit correction for multiple comparisons.

Our central question is: \textit{given a candidate phonetic key, does the resulting character sequence show statistically significant overlap with a known language corpus, beyond what random key assignment would produce?}

%------------------------------------------------------------
\section{Prior Work}
%------------------------------------------------------------

Scholarly literature falls into three categories. \textbf{Linguistic proposals} (Owens 1996; Achterberg et al.\ 2004; Weingarten 2016) assign phonetic values based on acrophonic principles or visual resemblance without statistical validation. \textbf{Structural analyses} (Sproat 2010; Rao et al.\ 2009) examine sign-distribution properties — entropy, Zipf distribution, positional bias — confirming the disc is likely a writing system without identifying the language. \textbf{Computational approaches} remain rare; no prior study applies Bonferroni-corrected Monte Carlo simulation across multiple competing keys against independently compiled reference corpora.

%------------------------------------------------------------
\section{Data}
%------------------------------------------------------------

\subsection{The Phaistos Disc}

The disc contains 45 distinct signs with 241 total occurrences in 61 word-groups (31 on Side~A, 30 on Side~B). Reading direction is outside-to-center (spiral inward) on both sides. Key structural observations, independent of any phonetic assumption:

\begin{itemize}
    \item Sign~\#2 is word-initial in 48\% of all words
    \item Sign~\#11 is word-final in 30\% of all words
    \item Shannon entropy $H = 3.045$ bits — consistent with syllabic writing (2.8--3.5~bits)
    \item Zipf $R^2 = 0.673$ — formulaic/ritual register
\end{itemize}

\subsection{Reference Corpora}

\paragraph{AED-TEI Egyptian Corpus.}
Akademie der Wissenschaften, Berlin: 675,773 tokens from 13,950 texts. Ritual subcorpus: 95,162 tokens from 1,370 texts (Pyramid Texts, Book of the Dead, Coffin Texts). License: CC-BY-SA~4.0.

\paragraph{Luwian Hieroglyphic Vocabulary.}
19 lexical entries with established phonetic values (Masson 1961; Hawkins 2000; Melchert 2003), independently attested in Anatolian inscriptions contemporary with the disc ($\sim$2000--1200~BCE).

\paragraph{Linear A Frequency Table.}
30 sign-syllable correspondences based on frequency-matched Linear A inventory. Linear A remains undeciphered; values are extrapolated from Linear~B cognates.

%------------------------------------------------------------
\section{Methodology}
%------------------------------------------------------------

\subsection{Blind Grid Test}

Ten phonetic keys (A--J) were constructed independently and evaluated simultaneously without post-hoc modification. Keys span Linear A acrophonic (\textsc{A\_Evans}), Linear A frequency (\textsc{B\_Freq}), Egyptian general corpus (\textsc{E1\_Egypt}), Egyptian Osiris-focused (\textsc{E2\_Wsir}), Cypriot syllabary (\textsc{F\_Cypriot}), Luwian Hieroglyphic (\textsc{G\_Luwian}), pure abjad (\textsc{H\_Abjad}), Linear A morphological (\textsc{I\_Morpho}), and Monte Carlo null (\textsc{J\_Null}).

\subsection{Scoring Function}

For key $K$ and word $W = [s_1, s_2, \ldots, s_n]$:
\begin{equation}
\text{score}(W, K) = \sum_{i=1}^{n} \max\_\text{length\_match}\!\left(K(s_i),\, \mathcal{V}\right)
\end{equation}
where $\max\_\text{length\_match}$ finds the longest vocabulary token $v \in \mathcal{V}$ matching the transliterated sequence at position $i$ without substring overlap. Total $S(K) = \sum_W \text{score}(W, K)$.

\subsection{Monte Carlo Null Distribution}

Key~J generates 10,000 random phonetic mappings (uniform random syllable assignment, seed~=~42), yielding:
\[
\mu_\text{null} = 151.9, \quad \sigma_\text{null} = 77.0
\]
Bonferroni threshold ($p < 0.005$, corrected for 10~keys): $S > 379$.
Publication-grade threshold ($p < 0.0001$): $S > 521$.

\subsection{Corpus-Domain Control}

G\_Luwian vocabulary was scored separately against theological and administrative AED-TEI subcorpora. A language-driven match should score similarly across registers; a register-driven match should score high only in the matching register.

\subsection{Negative Control}

1,000 synthetic discs were generated with: (a)~identical marginal sign frequencies; (b)~randomized sign adjacency. G\_Luwian was scored against each synthetic disc to isolate the sequential contribution from the frequency contribution.

%------------------------------------------------------------
\section{Results}
%------------------------------------------------------------

\subsection{Key Rankings}

\begin{table}[h]
\centering
\caption{Bonferroni-corrected key rankings. Threshold: $S > 379$ ($p < 0.005$).}
\begin{tabular}{llrrrrl}
\toprule
Rank & Key & Score & $Z$ & $p$-value & Bonf. & Refrain \\
\midrule
1 & \textsc{G\_Luwian} & \textbf{523} & 4.82 & \textbf{$<$0.0001} & $\checkmark\checkmark\checkmark$ & \textit{za-wa-tar} \\
2 & \textsc{E1\_Egypt}  & 491 & 4.40 & 0.0001  & $\checkmark\checkmark$ & \textit{n-m-r} \\
3 & \textsc{B\_Freq}    & 430 & 3.61 & 0.0009  & $\checkmark\checkmark$ & \textit{a-sa-ra} \\
4 & \textsc{I\_Morpho}  & 426 & 3.56 & 0.0009  & $\checkmark\checkmark$ & \textit{a-ku-te} \\
5 & \textsc{E2\_Wsir}   & 261 & 1.42 & 0.09    & ---  & \textit{A-sa-r} \\
6 & \textsc{A\_Evans}   & 188 & 0.47 & 0.30    & ---  & --- \\
7 & \textsc{F\_Cypriot} & 156 & 0.05 & 0.45    & ---  & \textit{a-ku-se} \\
8 & \textsc{H\_Abjad}   & 0   & $-1.97$ & 1.00 & ---  & \textsc{excluded} \\
\bottomrule
\end{tabular}
\end{table}

\textsc{H\_Abjad} scoring zero confirms the disc is not a pure consonantal script.

\subsection{Three Key-Independent Pillars}

\paragraph{Pillar~1 --- [Sign\#36$\to$Sign\#11] bigram ($p \approx 0$).}
Observed: 17 co-occurrences. Expected under independence: 2.2. Ratio: $7.69\times$. $Z = 10.00$. Under \textsc{G\_Luwian} this corresponds to \textit{wa-tar} (water, PIE~*\textit{wódr̥}). The excess cannot be explained by marginal frequencies alone.

\paragraph{Pillar~2 --- Corpus-domain control ($Z = 27.16$).}
\textsc{G\_Luwian} vocabulary scores $Z = 27.16$ against the theological AED-TEI subcorpus and $Z = -0.40$ against the administrative subcorpus. A $27.5\sigma$ gap between register categories confirms the disc belongs to a ritual register, independently of which phonetic key is correct.

\paragraph{Pillar~3 --- Sign~\#45 at centers only.}
Sign~\#45 (solar rosette) appears at exactly two locations: A31 (center of Side~A) and B30 (center of Side~B). This geometric observation is phonetically neutral. Probability of 14-occurrence sign landing exclusively at the two structural centers by chance: $p \approx 0.0018$.

\subsection{Center Chiasmus}

\begin{align*}
\text{A31} &= [45,\; 2,\; 36,\; 11,\; 22] \\
\text{B30} &= [45,\; 36,\; 11,\; 2,\; 6]
\end{align*}

Shared signs $\{2, 11, 36, 45\}$ appear in reversed index order, constituting a structural chiasmus consistent with intentional bilateral symmetric composition.

\subsection{Sensitivity and Cross-Validation}

\textbf{Sensitivity:} 105/105 single-pair perturbations of \textsc{G\_Luwian} remain above the Bonferroni threshold.

\textbf{Cross-validation:} Spearman $\rho = 0.779$ between Side~A and Side~B word-level scores. Transfer accuracy: A$\to$B = 64.6\%, B$\to$A = 91.6\%.

%------------------------------------------------------------
\section{Negative Control and Self-Critique}
%------------------------------------------------------------

\subsection{Frequency-Driven Token Score}

Mean G\_Luwian $Z$ on synthetic discs: $1.99$ (not significant at $p < 0.05$). Approximately 94\% of the token-level score is explained by marginal sign frequencies. \textbf{The token score $Z = 8.58$ (previously reported) is withdrawn as a primary argument and reclassified as exploratory.}

\subsection{Key Design Circularity}

\textsc{G\_Luwian} was constructed with awareness of disc statistics. This is the primary unresolved limitation. Remedy: blind replication by a Luwianologist with no prior knowledge of the key.

%------------------------------------------------------------
\section{Discussion}
%------------------------------------------------------------

\subsection{Primary Interpretation}

Under \textsc{G\_Luwian}: refrain $=$ \textit{za-wa-tar} (``this water'', PIE~*\textit{wódr̥}); center A31 $=$ \textit{ti-wa-za-wa-tar-ha} (``TIWAT! this water --- yes!''); center B30 $=$ \textit{ti-wa-wa-tar-za-an} (``TIWAT! water-judge --- here!''). The reading is consistent with a solar-water cosmological hymn: sun deity Tiwat descends into primordial waters (Side~A) and ascends reborn (Side~B).

\subsection{Egyptian Structural Parallel}

The \textit{Amduat} describes Ra's nightly spiral descent into Nun (primordial waters), union with Osiris at midnight, and solar rebirth. The structural mapping is: Ra~$\leftrightarrow$~Tiwat; Nun/Osiris~$\leftrightarrow$~\textit{watar}; midnight union~$\leftrightarrow$~disc centers A31/B30; descent~$\leftrightarrow$~Side~A; ascent~$\leftrightarrow$~Side~B. This parallel confirms descent/ascent solar theology but does not prove Luwian language.

%------------------------------------------------------------
\section{Limitations}
%------------------------------------------------------------

\begin{enumerate}
    \item \textbf{Key design circularity:} G\_Luwian constructed with awareness of disc statistics. Exploratory only until blind replication.
    \item \textbf{Hapax legomenon:} No second Phaistos-type text exists.
    \item \textbf{Token score:} $\sim$94\% frequency-driven; not a primary claim.
    \item \textbf{Vocabulary size:} G\_Luwian covers only 19 entries.
    \item \textbf{Linear A:} B\_Freq extrapolates from undeciphered script.
\end{enumerate}

%------------------------------------------------------------
\section{Conclusions}
%------------------------------------------------------------

We have demonstrated: (1)~a statistically non-random [Sign\#36$\to$Sign\#11] bigram ($Z=10$); (2)~identifiable ritual register (domain control $Z=27.16$ vs $Z=-0.40$); (3)~geometric chiasmus and Sign~\#45 exclusively at centers; (4)~\textsc{G\_Luwian} achieving $p<0.0001$ among 10~tested keys; (5)~a coherent solar-water cosmological reading with Amduat parallels; (6)~token scores are $\sim$94\% frequency-driven --- primary claims rest on key-independent evidence.

The blind multi-key grid methodology with Bonferroni correction, corpus-domain control, perturbation analysis, and negative control constitutes a replicable framework applicable to any undeciphered script with available reference corpora.

\textbf{Independent replication by a Luwianologist remains the critical next step.}

%------------------------------------------------------------
\begin{thebibliography}{99}
%------------------------------------------------------------
\bibitem{achterberg2004} Achterberg, W., Best, J., Enzler, K., Rietveld, L., \& Woudhuizen, F. (2004). \textit{The Phaistos Disc: A Luwian Letter to Nestor}. Dutch Monographs on Ancient History and Archaeology.

\bibitem{assmann2001} Assmann, J. (2001). \textit{The Search for God in Ancient Egypt}. Cornell University Press.

\bibitem{faulkner1969} Faulkner, R.O. (1969). \textit{The Ancient Egyptian Pyramid Texts}. Oxford University Press.

\bibitem{hawkins2000} Hawkins, J.D. (2000). \textit{Corpus of Hieroglyphic Luwian Inscriptions}. De Gruyter.

\bibitem{hornung1999} Hornung, E. (1999). \textit{The Ancient Egyptian Books of the Afterlife}. Cornell University Press.

\bibitem{masson1961} Masson, E. (1961). \textit{Recherches sur les plus anciens emprunts sémitiques en grec}. Paris.

\bibitem{melchert2003} Melchert, H.C. (2003). \textit{The Luwians}. Brill.

\bibitem{owens1996} Owens, G. (1996). The Phaistos Disc: A New Approach. \textit{Cretan Studies}, 5, 1--24.

\bibitem{rao2009} Rao, R.P.N., et al.\ (2009). Entropic Evidence for Linguistic Structure in the Indus Script. \textit{Science}, 324, 1165.

\bibitem{schweitzer2011} Schweitzer, S.D. (2011). AED-TEI Egyptian corpus. \url{https://github.com/simondschweitzer/aed-tei} (CC-BY-SA~4.0).

\bibitem{sproat2010} Sproat, R. (2010). Ancient Symbols, Computational Linguistics, and the Reviewing Practices of the General Science Journals. \textit{Computational Linguistics}, 36(3), 585--594.

\bibitem{weingarten2016} Weingarten, J. (2016). The Phaistos Disc: Pedigree of a Forgery. \textit{Journal of Prehistoric Religion}, 25.
\end{thebibliography}

\end{document}
